\begin{table*}[t]\centering\footnotesize
\setlength{\tabcolsep}{4pt}
\begin{tabular}{@{}lrrrrrrr@{}}
\toprule
Cue form & Health lit. & Confid. & Comm. & Emot. & Fluency & Median & Positive \\
\midrule
Author-written descriptor & 6.2 & 2.5 & 2.8 & 1.0 & 1.1 & 2.5 & 5/5 \\
Mechanical 2nd$\to$3rd person & 3.7 & 1.9 & 5.4 & 2.4 & 2.6 & 2.6 & 5/5 \\
Verbatim corpus instruction & 3.7$^\dagger$ & 2.5 & 15.5$^\dagger$ & 1.9 & 4.8$^\dagger$ & 3.7 & 5/5 \\
\bottomrule
\end{tabular}
\caption{\textbf{Cue-operationalisation robustness}, as style/placebo ratio. Unit of analysis is the dimension
($n=5$), exact sign test, $p=.0312$ for every arm --- the floor at this $n$. All three arms use one phrasing
per dimension, matched to the raw arm, which has exactly one instruction per (dimension, variant) in the
corpus; the eight-phrasing values in Table~\ref{tab:baseline} are the primary baseline.
$^\dagger$ clinical content words appear on one side of the cue only. Percentage magnitudes are \emph{not}
comparable across arms: raw instructions run 73--136 tokens against 11--15 for descriptors, which is why the
length-controlled ratio is tabulated. Spearman $\rho$ between descriptor and verbatim orderings
$= +0.50$ ($p=.39$, $n=5$).}
\label{tab:cue}
\end{table*}